\uuid{f0X9}
\titre{Bijection}
\auteur{bodin ; hanani ; mzari}
\organisation{exo7}
\licence{CC BY-NC-SA 4.0}
\datecreate{2020-01-20}
\niveau{L1}
\module{Analyse}
\chapitre{Continuité, limite et étude de fonctions réelles}
\sousChapitre{Etude de fonctions}
\difficulte{3}
\theme{Maximum, bijection, bijection}
\exo7qcmid{lille-1.continuite.030}

\contenu{
\begin{question}
\qclassification{123.04}
\qtags{motcle=bijection}
Soit $f : I \to J$ une fonction continue, où $I$ et $J$ sont des intervalles de $\Rr$. Quelles sont les assertions vraies ?
\begin{answers}
    \good{Si $f$ surjective et strictement croissante, alors $f$ est bijective.}

    \good{Si $f$ bijective, alors sa bijection réciproque $f^{-1}$ est continue.}

    \bad{Si $f$ bijective et $I=\Rr$, alors $J$ n'est pas un intervalle borné.}

    \good{Si $f$ bijective et $J$ est un intervalle fermé et borné, alors $I$ est un intervalle fermé et borné.}
\end{answers}
\begin{explanations}
La bijection réciproque d'une fonction continue est continue. En particulier cela implique que pour $f^{-1} : J \to I$, si $J$ est un intervalle fermé et borné, alors $I$ aussi.
\end{explanations}
\end{question}
}
